% Tailored from reusable AI/ML resume for CyberCoders Senior Machine Learning Engineer.
\documentclass[9pt,letterpaper]{extarticle}
\usepackage[left=0.68in,right=0.68in,top=0.45in,bottom=0.45in]{geometry}
\usepackage[T1]{fontenc}
\input{glyphtounicode}
\pdfgentounicode=1
\usepackage[utf8]{inputenc}
\usepackage{lmodern}
\usepackage{microtype}
\usepackage{pifont}
\usepackage{xcolor}
\usepackage{enumitem}
\usepackage{titlesec}
\usepackage[hidelinks]{hyperref}
\usepackage{parskip,ragged2e,needspace}
\pagestyle{empty}
\setlength{\parindent}{0pt}
\setlength{\parskip}{1pt}
\definecolor{accent}{HTML}{155B8A}
\titleformat{\section}{\large\bfseries\color{accent}}{}{0pt}{}[\vspace{0.5pt}\color{accent}\titlerule]
\titlespacing*{\section}{0pt}{3pt}{2pt}
\setlist[itemize]{leftmargin=1.15em,itemsep=0.7pt,topsep=1pt,parsep=0pt,partopsep=0pt}
\newcommand{\role}[4]{\Needspace{4\baselineskip}\textbf{#1}\hfill\textbf{#2}\\[-1pt]\textit{#3}\hfill\textit{#4}\par}
\newcommand{\project}[3]{\Needspace{5\baselineskip}\textbf{#1}\hfill\textit{#2 \textbar{} #3}\par}
\newcommand{\cert}[3]{\Needspace{4\baselineskip}\textbf{#1}\hfill\textit{#2}\par{\small #3\par}}
\newcommand{\pub}[3]{\Needspace{3\baselineskip}\textbf{#1}\hfill\textit{#2}\par{\small #3\par}}
\hypersetup{
  pdftitle={Waseem Raza, PhD - Senior Machine Learning Engineer},
  pdfauthor={Waseem Raza}
}

\begin{document}
\begin{center}
{\LARGE\bfseries Waseem Raza, PhD}\\[2pt]
{\large\bfseries Senior Machine Learning Engineer | Generative Modeling \& Production ML}\\[1pt]
{\small Python | Deep Learning | GANs \& Autoencoders | Data Pipelines | MLOps}\\[3pt]
{\footnotesize New Jersey, United States \quad
\ding{37}\enspace \href{tel:+17328039703}{+1 732-803-9703} \quad
\ding{41}\enspace \href{mailto:waseem.placements@gmail.com}{waseem.placements@gmail.com} \textbar{}
\href{mailto:waseemraza844@gmail.com}{waseemraza844@gmail.com}\\[1pt]
\href{https://www.linkedin.com/in/waseem-raza-phd-9a737858/}{LinkedIn} \textbar{}
\href{https://github.com/WaseemRaza844}{GitHub} \textbar{}
\href{https://scholar.google.com/citations?user=d58Drr0AAAAJ&hl=en}{Google Scholar} \textbar{}
\href{https://waseemraza844.github.io/portfolio/}{Portfolio}}
\end{center}
\vspace{-5pt}

\section{Professional Summary}
Senior-level Machine Learning Engineer and researcher with a PhD in Electrical and Computer Engineering and 7+ years of combined industry and research experience designing, training, validating, and operationalizing ML systems. Strong Python and deep-learning background spanning PyTorch, TensorFlow/Keras, scikit-learn, XGBoost, NumPy, SciPy, Pandas, GANs, autoencoders, transfer learning, domain adaptation, reinforcement learning, CNN/LSTM models, anomaly detection, and predictive modeling. Experience combines research-quality experimentation with production-oriented engineering: reusable data/model pipelines, Git-based development, CI/CD, Docker/Kubernetes, cross-domain validation, cloud analytics, and stakeholder-driven deployment. Led 100+ controlled ML experiments, analyzed 500K+ operational logs, improved autoencoder-based compression by 25\%, and delivered production-oriented ML workflows across AT\&T, Futurewei, university research, and geoscience analytics.

\section{Technical Skills}
\textbf{Generative Modeling \& Representation Learning:} GANs, autoencoders, deep neural networks, synthetic-data generation, transfer learning, domain adaptation, representation learning, transformer fine-tuning concepts, generative-model evaluation.\\
\textbf{ML / Deep Learning:} PyTorch, TensorFlow, Keras, scikit-learn, XGBoost, Hugging Face Transformers; CNNs, LSTMs, reinforcement learning, classification, regression, anomaly detection, time-series modeling, model evaluation.\\
\textbf{Python \& Numerical Computing:} Python, NumPy, SciPy, Pandas, SQL, C/C++, Jupyter, numerical experimentation, statistical analysis, vectorized processing, Linux, Bash/CLI.\\
\textbf{Training, Inference \& Performance:} preprocessing, feature engineering, train/validation/test design, hyperparameter experimentation, cross-validation, batch inference, metric tracking, robustness testing, failure analysis, distribution-shift evaluation, model comparison.\\
\textbf{Data Pipelines \& Large-Scale Processing:} Azure Databricks, Apache Spark, PySpark, Synapse Analytics, Azure Functions, ETL/ELT, structured/unstructured data, data-quality validation, AI-ready datasets, synthetic-data pipelines.\\
\textbf{MLOps \& Software Engineering:} Git, GitHub Actions, CI/CD, Docker, Kubernetes, OpenShift, testing, REST APIs, Flask, Linux, reproducible repositories, deployment-oriented model validation, technical documentation.\\
\textbf{Collaboration \& Leadership:} technical ownership, cross-functional requirements translation, research mentoring, experiment review, reproducibility standards, technical writing, presentations, and knowledge transfer.

\section{Professional Experience}
\role{Machine Learning Engineer}{02/2025 -- Present}{AT\&T, contractor via Innovatix Inc.}{Bedminster, NJ}
\begin{itemize}
\item Own end-to-end development, testing, and operational integration of ML diagnostics for critical E911 service-assurance workflows, improving automated root-cause classification accuracy from 20\% to 80\%.
\item Design domain-adaptive neural models and reusable data/model pipelines that improve generalization across heterogeneous operational domains and distribution shifts.
\item Develop multimodal AI prototypes combining text, location, network context, and operational signals; build error-analysis and failure-mode assessment procedures for deployment readiness.
\item Partner with network engineers, software teams, and public-safety stakeholders to translate operational objectives into measurable ML requirements, validated datasets, model-selection decisions, and deployable workflows.
\end{itemize}

\role{Data Scientist and AI/ML Engineer}{09/2024 -- 12/2024}{Impac Exploration Services}{Weatherford, OK}
\begin{itemize}
\item Built ML and LLM-assisted analytical workflows for large geochemical datasets, collaborating with geologists to translate scientific objectives into scalable data and modeling pipelines.
\item Developed ingestion, cleansing, feature-preparation, transformation, and validation workflows using Python, Azure Databricks, Azure Functions, and Synapse Analytics.
\item Automated ETL processing, reducing data-preparation time by 40\%; implemented cross-validation and data-quality checks that improved reliability of analytical outputs by 20\%.
\end{itemize}

\newpage
\role{Graduate Research Assistant}{09/2019 -- 07/2025}{AI4Networks, University of Oklahoma}{Tulsa, OK}
\begin{itemize}
\item Led research spanning generative modeling, robust ML, reinforcement learning, domain adaptation, wireless digital twins, network optimization, and wearable intelligence.
\item Built transfer-learning, adversarial-learning, GAN, autoencoder, CNN/LSTM, and reinforcement-learning pipelines for sparse-data, cross-domain, and sequence-modeling problems.
\item Designed reusable Python simulation/data-generation systems and robustness-evaluation frameworks; led 100+ controlled experiments from formulation through quantitative validation and publication.
\item Mentored junior researchers in ML implementation, experiment design, reproducibility, technical writing, and presentation of results.
\end{itemize}

\role{Advanced Wireless Technology Development Intern}{09/2022 -- 01/2023}{Futurewei Technologies}{Schaumburg, IL}
\begin{itemize}
\item Improved autoencoder-based channel-state-information compression by 25\% using transfer learning and systematic experimentation on large wireless datasets.
\item Conducted 50+ experiments to evaluate reconstruction performance, robustness, and generalization across changing data conditions; translated results into recommendations for next-generation communication systems.
\end{itemize}

\role{Advanced-Degrees Intern}{06/2024 -- 08/2024}{AT\&T Labs}{Austin, TX}
\begin{itemize}
\item Analyzed 500K+ LTE/5G-NR diagnostic logs and applied ML methods to detect anomalies, reducing diagnostic time by 20\%.
\item Developed and validated data-driven decision logic through structured experimentation, performance analysis, and communication of results to engineering teams.
\end{itemize}

\section{Education}
\textbf{PhD, Electrical and Computer Engineering}\hfill\textit{09/2019 -- 07/2025}\\
University of Oklahoma, Norman, Oklahoma, USA\\
\textbf{MSc, Telecommunication Engineering}\hfill\textit{08/2014 -- 08/2016}\\
University of Engineering and Technology, Taxila, Pakistan\\
\textbf{BSc, Telecommunication Engineering}\hfill\textit{11/2010 -- 06/2014}\\
University of Engineering and Technology, Taxila, Pakistan

\section{Selected Machine Learning Projects}
\project{Domain-Informed GANs for Propagation Models \& Digital Twins}{Research Project}{2024}
Developed generative-modeling workflows for wireless digital twins using 3GPP-informed neural pretraining, GAN-based synthetic-data generation, and scarce measured data. Compared synthetic-data quality and downstream regression performance, and evaluated robustness under domain shift and limited real training data.\\[-1pt]
{\small\textit{Tools \& methods:} Python, GANs, deep neural networks, synthetic data, domain adaptation, transfer learning, regression, RMSE/R\textsuperscript{2}, statistical validation, distribution-shift evaluation.}

\project{Autoencoder CSI Compression \& Transfer Learning}{Industry Research Project}{2022--2023}
Developed and evaluated autoencoder-based representation/compression workflows for large channel-state-information datasets. Applied transfer learning to improve performance by 25\%, executed 50+ controlled experiments, and benchmarked reconstruction accuracy and generalization across changing data conditions.\\[-1pt]
{\small\textit{Tools \& methods:} Python, autoencoders, deep learning, representation learning, transfer learning, signal compression, train/validation experimentation, performance benchmarking, robustness analysis.}

\project{Wearable Intelligence for Health Monitoring}{Research Project}{2024--2026}
Developed deep-learning pipelines for medication-taking and dietary-activity recognition from smartwatch accelerometer/gyroscope data. Implemented segmentation, feature engineering, CNN/LSTM-style sequence models, cross-user validation, and recall-oriented error analysis for real-world health-monitoring scenarios.\\[-1pt]
{\small\textit{Tools \& methods:} Python, CNN/LSTM, time-series modeling, sensor preprocessing, feature engineering, sequence classification, cross-user evaluation, confusion matrices, model comparison.}

\project{SyntheticNET: 3GPP-Compliant Network Simulation}{Research Project}{2020--2024}
Built a reusable Python simulation and synthetic-data environment for large-scale experiments where measured data is sparse or expensive. Generated controlled mobility, propagation, and network-state datasets to support repeatable model training, benchmarking, and robustness studies.\\[-1pt]
{\small\textit{Tools \& methods:} Python, numerical simulation, synthetic data, experiment automation, reproducibility, 3GPP models, data pipelines, performance evaluation.}

\project{AI-Ready Lakehouse Pipeline}{Portfolio Project}{2024}
Designed scalable ingestion, transformation, validation, and feature-preparation workflows for analytics and downstream ML using a cloud lakehouse architecture. Focused on repeatability, data quality, and scalable processing for model development and inference workflows.\\[-1pt]
{\small\textit{Tools \& methods:} Azure Databricks, Apache Spark, PySpark, ETL/ELT, data ingestion, schema transformation, data validation, feature preparation, cloud data engineering.}

\project{Cloud Computer Vision}{Guided Project}{2025}
Implemented image-classification, object-detection, and vehicle-counting workflows using managed cloud models and Python inference pipelines. Validated preprocessing, prediction outputs, class/count aggregation, and model behavior across visual inputs.\\[-1pt]
{\small\textit{Tools \& methods:} Python, TensorFlow concepts, computer vision, object detection, image classification, preprocessing, inference, model evaluation, Google Cloud AutoML.}

\section{Production ML, Reproducibility \& Technical Leadership}
\begin{itemize}
\item Build ML workflows as reusable stages for data preparation, model training/inference, metric computation, error analysis, and validation; document assumptions, experiments, and acceptance criteria.
\item Use Git, automated tests, GitHub Actions, CI/CD, Docker, Kubernetes, and command-line workflows to support reproducibility and production-oriented engineering practices.
\item Review model behavior for data leakage, metric misuse, shape/type errors, cross-domain failures, and mismatches between intended and actual model behavior before deployment decisions.
\item Mentor junior researchers and collaborate with engineers, scientists, and domain experts to translate research objectives into reproducible implementations and measurable results.
\end{itemize}

\section{Selected Technical Training}
\cert{\href{https://www.coursera.org/account/accomplishments/verify/VU3Y5TBLVLZT}{Machine Learning}}{Stanford Online / Coursera | Completed 03/2020}{Supervised/unsupervised learning, regression, classification, regularization, clustering, bias/variance, optimization concepts, and practical model evaluation.}
\cert{\href{https://www.coursera.org/account/accomplishments/verify/SG4GDJ4CLPVL}{Generative AI Engineering and Fine-Tuning Transformers}}{IBM / Coursera | Completed 08/2026}{Transformer architectures, pretrained models, model adaptation, fine-tuning strategies, tokenization concepts, training/inference tradeoffs, and LLM evaluation.}
\cert{\href{https://www.coursera.org/account/accomplishments/verify/Y6NP4GWLS4DZ}{Introduction to Deep Learning \& Neural Networks with Keras}}{IBM / Coursera | Completed 08/2026}{Neural-network architecture, activation functions, optimization, model training, validation, and iterative performance improvement.}
\cert{\href{https://www.coursera.org/account/accomplishments/verify/R6248BNQAC0D}{Introduction to Computer Vision with TensorFlow}}{Google Cloud / Coursera | Completed 08/2024}{Image preprocessing, convolutional-network concepts, image classification, model construction, training, validation, and evaluation.}

\section{Selected Research Evidence}
\pub{\href{https://doi.org/10.1109/PIMRC59610.2024.10817403}{An AI-Driven Framework for Enhancing Resilience in Propagation Models to Enable Digital Twin}}{IEEE PIMRC, 2024}{Generative augmentation, neural modeling, scarce-data learning, robustness evaluation, and downstream predictive-performance validation.}
\pub{Real-World Medication Adherence Monitoring Using Hybrid Neural Networks and Smartwatches}{HealthINF, Accepted 2026}{Deep-learning sequence classification, multimodal sensor data, cross-user validation, and real-world reliability analysis.}
\pub{Data-Driven Intelligent Outage Management for High Shadowing Environments in 5G\&B Networks}{IEEE SmartNets, 2024}{Machine-learning-based anomaly/outage detection, model evaluation, and decision support in challenging operational environments.}
\pub{\href{https://doi.org/10.3390/s25247651}{On Multi-Parameter Optimization and Proactive Reliability in 5G and Beyond Cellular Networks}}{Sensors, 2025}{Predictive modeling, optimization, reliability analysis, and systematic evaluation for data-driven network operation.}

\end{document}
